% 06-section5.tex -- §5 重写：pattern-bound Roofline（相图降级为 5.5）
\section{Boundary: a pattern-bound roofline}\label{sec:roofline}

The cost of a sparse factorization is set by its \emph{pattern}, not by $n$: the
fill determines both the arithmetic and the traffic, so the usual roofline
x-axis is a function of the elimination order. We make that dependence explicit
and use it to decide what the remaining levers are.

\subsection{Arithmetic intensity of a factorization}

For a factorization with $\mathrm{nnz}(L+U)=f$ nonzeros, the dominant work is
the rank-$b$ updates; counting $8$ bytes per factor entry and $4$ per index,
\begin{equation}
  \mathrm{AI}(A,\mathcal{O}) \;=\; \frac{\mathrm{flops}(A,\mathcal{O})}{12\,f(A,\mathcal{O})},
  \qquad f = f(A,\mathcal{O}),\label{eq:ai}
\end{equation}
where $\mathcal{O}$ is the ordering. Because $f$ is the numerator of the fill
ratio and appears in the denominator of the intensity, \emph{choosing an ordering
moves the operating point horizontally}: $\mathrm{AI}\propto 1/f$ at fixed work.
This is the ``pattern-bound'' part: the roofline position is a property of the
elimination, not of the matrix alone.

\subsection{Where the points land on this device}

On the Mars X201 we measure $B=1.45$\,TB/s with a stream kernel and
$P=29.03$\,TFLOP/s with the vendor dense GEMM (\texttt{hcBLAS}, $n=4096$), giving
a ridge point of
\begin{equation}
  \mathrm{AI}^\ast \;=\; P/B \;=\; 20\ \text{FLOP/byte}.
\end{equation}
For the thirteen circuit matrices the intensity computed from \eqref{eq:ai} is
$10^{3}$--$10^{4}$ FLOP/byte, \emph{two orders of magnitude to the right of the
ridge}. Every matrix is therefore squarely compute-bound, and the bandwidth roof
plays no role at this problem size.

\subsection{Precision is not a speed lever here}

A compute-bound regime is exactly where lower precision is supposed to pay---but
only if the low-precision arithmetic is actually faster. On this device it is
not:
\begin{center}
\small
\begin{tabular}{lrr}
  \toprule
  kernel & throughput & ratio to FP64 \\
  \midrule
  \texttt{hcBLAS} DGEMM & $28.79$\,TFLOP/s & $1.000$ \\
  \texttt{hcBLAS} SGEMM & $29.03$\,TFLOP/s & $1.008$ \\
  our naive shared-memory GEMM, FP64 & $1.24$\,TFLOP/s & $0.043$ \\
  our naive shared-memory GEMM, FP32 & $2.82$\,TFLOP/s & $0.098$ \\
  \bottomrule
\end{tabular}
\end{center}
At vendor efficiency $s_{32}=1.008$: there is \emph{no} FP32-over-FP64 compute
advantage. The factor-of-$2.29$ that a naive kernel suggests is an artifact of
that kernel being limited to $875$\,GB/s of element throughput in FP32 while its
FP64 variant reaches $1446$\,GB/s---a bandwidth effect, not a compute one.
Consequently lowering the arithmetic precision cannot shorten the factorization
on this device; only a genuinely faster arithmetic path (tensor cores) or a
better kernel can.

\subsection{The cost model and its corner solution}\label{sec:cost}

Hold the algorithm fixed and let $s_u$ be the throughput of precision $u$
relative to FP64 and $k(\kappa)$ the refinement count. In FP64-equivalent units
the cost of a solve is
\begin{equation}\label{eq:cost}
  T(u) \;=\; \frac{1}{s_u}\,\frac{2n^{3}}{3} \;+\; 4k(\kappa)\,n^{2},
\end{equation}
the first term being the factorization and the second the refinement. Both
summands are monotone decreasing in $s_u$ over the feasible set, and the feasible
set is monotone in each slot, so the optimum is attained at a \emph{corner}: one
uses the fastest admissible precision for the factorization and the highest
admissible for the residual.

\begin{proposition}[Corner solution and crossover]\label{prop:corner}
Between two precisions with ratios $s$ and $s'$ and refinement counts $k<k'$,
$T$ is minimised by the cheaper per-flop option up to
\begin{equation}
  n^{\ast}(\kappa) \;=\; \frac{6\,(k'-k)\,s\,s'}{s'-s},
\end{equation}
beyond which the faster arithmetic wins despite the extra iterations.
\end{proposition}

\begin{proof}
Set $T(s)=T(s')$ in \eqref{eq:cost} and solve for $n$; the quadratic terms cancel
and the linear system in $n$ gives the stated root. Monotonicity of the feasible
set gives the corner property.
\end{proof}

\begin{remark}
Equation~\eqref{eq:cost} is the reason the feasibility question of
Section~\ref{sec:bounds} and the boundary question of this section are separate:
the first two terms of \eqref{eq:cost} are unaffected by whether the refinement
converges, and the convergence condition of Theorem~\ref{thm:dichotomy} involves
only $\kappa$ and $u_q$.
\end{remark}

\subsection{When a low-precision tier would pay}\label{par:impl}

The analysis above says that at $s_{32}=1.008$ the arithmetic precision is not a
lever. A low-precision tier becomes one only if its throughput exceeds a
threshold. Writing $\Delta k=k_{11}-k_{32}$, the 11-bit tier beats FP32 if and
only if
\begin{equation}
  s_{11} > s_{11}^{\ast}(\kappa,n) \;=\; \frac{1}{1/s_{32}-6\Delta k/n},
  \qquad\text{solvable only when } n>6\,s_{32}\,\Delta k .
\end{equation}
With the measured $s_{32}=1.008$ the thresholds collapse to $1.03$--$1.20$
(Table~\ref{tab:s11}): the tier pays as soon as it is marginally faster than
FP64, and the binding constraint is not the threshold but whether such a path
exists on the device at all.

\begin{table}[t]
  \centering
  \small
  \caption{Required $s_{11}$ for the 11-bit tier to pay, with $s_{32}=1.008$.}
  \label{tab:s11}
  \begin{tabular}{lrrrr}
    \toprule
    $\kappa$ & $\Delta k$ & $n>6s_{32}\Delta k$ & $s_{11}^\ast(4096)$ & $s_{11}^\ast(16384)$ \\
    \midrule
    $10^{6}$  & 79  & 477 & 1.14 & 1.04 \\
    $10^{8}$  & 110 & 665 & 1.20 & 1.05 \\
    $10^{10}$ & 78  & 471 & 1.14 & 1.04 \\
    $10^{12}$ & 48  & 290 & 1.08 & 1.03 \\
    \bottomrule
  \end{tabular}
\end{table}

\begin{figure}[t]
  \centering
  \includegraphics[width=\linewidth]{figures/fig10_s11_threshold.pdf}
  \caption{The implementation threshold $s_{11}^\ast(\kappa,n)$ at the measured
  $s_{32}=1.008$ (grey). A tier pays when its throughput lies above the curves.}
  \label{fig:s11}
\end{figure}

\subsection{The \texorpdfstring{\((\kappa,s)\)}{(kappa,s)} phase diagram}\label{sec:phase}

For completeness we record the picture that motivated the cost model: if both
ratios were free parameters, the optimal tier as a function of $(\kappa,s)$
consists of regions separated by $n^{\ast}$ contours
(Figure~\ref{fig:phase}). On the measured device the 11-bit region is empty
because $s_{11}\approx s_{32}$; the diagram is retained as the map of what a
device with a faster low-precision path would unlock.

\begin{figure}[t]
  \centering
  \includegraphics[width=0.85\linewidth]{figures/fig9_phase_n4096.pdf}
  \caption{Optimal tier as a function of $(\kappa,s)$ at $n=4096$, with
  $s_{11}=2.34$ as measured. The FP32 and 11-bit regions are separated by the
  $n^\ast$ contour of Proposition~\ref{prop:corner}.}
  \label{fig:phase}
\end{figure}
